% =====================================================================
%  CV bilíngue (single source) — Marcos Junior Bueno Selzler
%  Idioma PRIMÁRIO: Inglês (default ao compilar main.tex)
%  Idioma SECUNDÁRIO: Português  -> compile main-pt.tex
%
%  Como gerar cada PDF:
%    EN (default): compile este main.tex
%    PT          : compile main-pt.tex  (faz \def\CVPT{}% =====================================================================
%  CV bilíngue (single source) — Marcos Junior Bueno Selzler
%  Idioma PRIMÁRIO: Inglês (default ao compilar main.tex)
%  Idioma SECUNDÁRIO: Português  -> compile main-pt.tex
%
%  Como gerar cada PDF:
%    EN (default): compile este main.tex
%    PT          : compile main-pt.tex  (faz \def\CVPT{}% =====================================================================
%  CV bilíngue (single source) — Marcos Junior Bueno Selzler
%  Idioma PRIMÁRIO: Inglês (default ao compilar main.tex)
%  Idioma SECUNDÁRIO: Português  -> compile main-pt.tex
%
%  Como gerar cada PDF:
%    EN (default): compile este main.tex
%    PT          : compile main-pt.tex  (faz \def\CVPT{}% =====================================================================
%  CV bilíngue (single source) — Marcos Junior Bueno Selzler
%  Idioma PRIMÁRIO: Inglês (default ao compilar main.tex)
%  Idioma SECUNDÁRIO: Português  -> compile main-pt.tex
%
%  Como gerar cada PDF:
%    EN (default): compile este main.tex
%    PT          : compile main-pt.tex  (faz \def\CVPT{}\input{main.tex})
%  Macro de tradução:  \tr{<texto EN>}{<texto PT>}
% =====================================================================

\documentclass[a4paper,10pt]{article}
\usepackage[margin=0.5in,nofoot]{geometry}
\usepackage{fontawesome5}
\usepackage{xcolor}

\definecolor{accent}{HTML}{1F4E79}

% ---- Fonte (XeLaTeX/LuaLaTeX) ----
\usepackage{fontspec}
\setmainfont{Montserrat}[
    Path = ./fonts/,
    Extension = .ttf,
    UprightFont = *-Regular,
    BoldFont = *-Bold,
    ItalicFont = *-Italic
]

% ---- Idioma / hifenização ----
\usepackage{polyglossia}
\ifdefined\CVPT
  \setdefaultlanguage{portuguese}\setotherlanguage{english}
\else
  \setdefaultlanguage{english}\setotherlanguage{portuguese}
\fi

% ---- Macro de tradução: \tr{EN}{PT} (default = EN) ----
\newcommand{\tr}[2]{\ifdefined\CVPT #2\else #1\fi}

\usepackage{titlesec}
\usepackage{enumitem}

\usepackage{hyperref}
\hypersetup{
    colorlinks=true,
    linkcolor=accent,
    filecolor=accent,
    urlcolor=accent,
    citecolor=accent
}

\titleformat{\section}{\large\bfseries}{\thesection}{1em}{}[\titlerule]
\titlespacing*{\section}{0pt}{*1}{*1}

\newcommand{\entry}[4]{
  \noindent\textbf{#1} \hfill #2 \\
  \noindent\textit{#3} \hfill \textit{#4} \\
  \vspace{2pt}
}

\newcommand{\project}[2]{
  \noindent\textbf{#1} \hfill #2 \\
  \vspace{2pt}
}

\begin{document}
\small
\setlength{\parskip}{4pt}
\setlength{\parindent}{0pt}

\pagenumbering{gobble}

\noindent
\begin{minipage}[t]{0.5\textwidth}
{\Large \textbf{Marcos Junior Bueno Selzler}}

\vspace{0.4em}

\end{minipage}%
\begin{minipage}[t]{0.5\textwidth}
\raggedleft

\tr{Brazil}{Brasil} (GMT-3)

\href{tel:+5543988431018}{\faPhone \space (43) 98843-1018} \\
\href{mailto:h0wzymarcos@gmail.com}{\faEnvelope \space h0wzymarcos@gmail.com} \\
\href{https://www.linkedin.com/in/marcosh0wzy/}{\faLinkedin \space marcosh0wzy} \\
\href{https://github.com/H0wZy}{\faGithub \space H0wZy} \\
\href{https://linktr.ee/h0wzymarcos}{\faLink \space \tr{Portfolio}{Portfólio}}

\end{minipage}

\vspace{1em}

\begin{center}
    \textbf{SOFTWARE ENGINEER \textbar\ BACKEND \& FULL STACK}
\end{center}

\vspace{0.4em}

\noindent
\tr{%
Software Engineer at TCS (Copel account), focused on backend with Java/JSF and .NET/Golang. Two-time TCS AI Friday champion for solutions built with LLMs and predictive models. Experienced in modernizing legacy systems to the cloud (GCP/AWS) and building REST APIs with Clean Code, SOLID and testing.%
}{%
Engenheiro de Software na TCS (conta Copel), com foco em backend Java/JSF e .NET/Golang. Bicampeão do TCS AI Friday por soluções com LLMs e modelos preditivos. Experiência em modernização de sistemas legados para cloud (GCP/AWS) e construção de APIs REST com Clean Code, SOLID e testes.%
}

\section*{\tr{Professional Experience}{Experiência Profissional}}

\entry{Tata Consultancy Services \textbar\ Copel}{\faCalendar \space 05/2025 -- \tr{Present}{Presente}}{Software Engineer}{Londrina - PR}
\vspace{-0.4em}
\begin{itemize}[leftmargin=*, itemsep=2pt]
\setlength\itemsep{0em}
\item \tr{I develop and maintain Java 8 / JSF applications on WildFly/JBoss servers.}{Desenvolvo e mantenho aplicações Java 8 com JSF em servidores WildFly/JBoss.}
\item \tr{Optimized critical SQL queries, improving performance and stability.}{Otimizei queries SQL críticas, melhorando performance e estabilidade.}
\item \tr{Migrated applications from JBoss 7 to WildFly 23.}{Migrei aplicações de JBoss 7 para WildFly 23.}
\item \tr{Modernize on-premise legacy systems to GCP, with integrations via Denodo.}{Modernizo sistemas legados on-premise para o GCP, com integrações via Denodo.}
\item \tr{Build backend services in Golang.}{Desenvolvo serviços de backend em Golang.}
\item \tr{Work across cloud environments (GCP and AWS).}{Atuo em ambientes cloud (GCP e AWS).}
\end{itemize}

\entry{Tata Consultancy Services \textbar\ Bootcamp 2x2y}{\faCalendar \space 05/2025 -- 11/2025}{Bootcamper Software Engineer}{Londrina - PR}
\vspace{-0.4em}
\begin{itemize}[leftmargin=*, itemsep=2pt]
\setlength\itemsep{0em}
\item \tr{Intensive full-stack software engineering training at TCS.}{Formação intensiva em engenharia de software full stack pela TCS.}
\item \tr{Built hands-on projects applying layered architecture and Clean Code.}{Desenvolvi projetos práticos aplicando arquitetura em camadas e Clean Code.}
\end{itemize}

\section*{\tr{Awards \& Highlights}{Prêmios e Destaques}}

\entry{\tr{Two-time TCS AI Friday Champion}{Bicampeão TCS AI Friday}}{05/2026}{\tr{Pipeverse Project}{Projeto Pipeverse}}{\href{https://www.linkedin.com/posts/marcosh0wzy_tcs-tcsai-aifriday-activity-7463971325016006657-WzoZ}{\faLink \space LinkedIn}}
\vspace{-0.4em}
\begin{itemize}[leftmargin=*, itemsep=2pt]
\setlength\itemsep{0em}
\item \tr{Predictive system to estimate the probability of failure in AI projects (data quality, problem definition, adoption and monitoring).}{Sistema preditivo para estimar a probabilidade de falha em projetos de IA (qualidade de dados, definição do problema, adoção e monitoramento).}
\end{itemize}

\entry{\tr{TCS AI Friday Champion}{Campeão TCS AI Friday}}{02/2026}{\tr{Energy Saver Project}{Projeto Energy Saver}}{\href{https://www.linkedin.com/posts/marcosh0wzy_ai-llm-dataanalysis-activity-7439643718057385985--a7G}{\faLink \space LinkedIn}}
\vspace{-0.4em}
\begin{itemize}[leftmargin=*, itemsep=2pt]
\setlength\itemsep{0em}
\item \tr{LLM-based solution for energy-consumption analysis from CSV data: anomaly detection and automated recommendations to reduce waste.}{Solução com LLM para análise de consumo energético a partir de CSV: detecção de anomalias e recomendações automatizadas para redução de desperdício.}
\end{itemize}

\section*{\tr{Skills}{Habilidades}}

\begin{itemize}[leftmargin=*, itemsep=2pt]
\setlength\itemsep{0em}
\item \textbf{Tech Stack}: .NET, C\#, Golang, Java, Spring Boot, React, Next.js, TypeScript, SQL/NoSQL, Docker
\item \textbf{Cloud \& DevOps}: GCP (Cloud Run, Cloud SQL, GCS), AWS, Terraform (IaC), Docker, CI/CD, Git/GitHub
\item \textbf{\tr{AI \& LLMs}{IA \& LLMs}}: \tr{data analysis with LLMs, predictive models, prompt engineering}{análise de dados com LLMs, modelos preditivos, prompt engineering}
\item \textbf{\tr{Concepts}{Conceitos}}: \tr{Clean Code, SOLID, DDD, OOP, REST APIs, Observability, Unit Testing (xUnit, Moq)}{Clean Code, SOLID, DDD, POO, APIs REST, Observabilidade, Testes Unitários (xUnit, Moq)}
\item \textbf{\tr{Languages}{Idiomas}}: \tr{English (Intermediate), Portuguese (Native)}{Inglês (Intermediário), Português (Nativo)}
\end{itemize}

\section*{\tr{Projects}{Projetos}}

\project{Telas Paraná --- \tr{Institutional Website}{Site Institucional}}{\href{https://frontend-4mxyivrfva-r.a.run.app/}{\faLink \space \tr{Live}{Produção}} \textbar\ \href{https://telasearamesparana.com.br}{\tr{previous site}{site anterior}}}
\vspace{-0.4em}
\begin{itemize}[leftmargin=*, itemsep=2pt]
\setlength\itemsep{0em}
\item \tr{Full-stack redesign of the company's institutional website, replacing the legacy site.}{Redesign full-stack do site institucional da empresa, substituindo o site legado.}
\item \tr{React / Next.js front-end and two .NET/C\# microservices (Users.Api, Leads.Api), all deployed on GCP Cloud Run.}{Front-end em React / Next.js e dois microsserviços .NET/C\# (Users.Api, Leads.Api), todos implantados no GCP Cloud Run.}
\item \tr{Cloud SQL database, Cloud Storage (GCS) for media, and infrastructure provisioned as code with Terraform.}{Banco no Cloud SQL, mídia no Cloud Storage (GCS) e infraestrutura provisionada como código com Terraform.}
\end{itemize}

\project{C\# User API}{\href{https://github.com/H0wZy/UserApi}{\faLink \space GitHub}}
\vspace{-0.4em}
\begin{itemize}[leftmargin=*, itemsep=2pt]
\setlength\itemsep{0em}
\item \tr{REST API in C\#/.NET with Entity Framework and PostgreSQL, layered architecture and user CRUD.}{API REST em C\#/.NET com Entity Framework e PostgreSQL, arquitetura em camadas e CRUD de usuários.}
\item \tr{JWT authentication, authorization control and test coverage with xUnit and Moq.}{Autenticação JWT, controle de autorização e cobertura de testes com xUnit e Moq.}
\end{itemize}

\project{Golang User API}{\href{https://github.com/H0wZy/user-api}{\faLink \space GitHub}}
\vspace{-0.4em}
\begin{itemize}[leftmargin=*, itemsep=2pt]
\setlength\itemsep{0em}
\item \tr{REST API in Golang with Gin and Gorm: user CRUD, authentication and authorization.}{API REST em Golang com Gin e Gorm: CRUD de usuários, autenticação e autorização.}
\item \tr{Backend structured with a focus on performance and code organization.}{Backend estruturado com foco em performance e organização de código.}
\end{itemize}

\project{Web EasyDonate}{\href{https://github.com/H0wZy/web.easydonate}{\faLink \space GitHub}}
\vspace{-0.4em}
\begin{itemize}[leftmargin=*, itemsep=2pt]
\setlength\itemsep{0em}
\item \tr{Donation marketplace in .NET with PostgreSQL, oriented to circular economy and scalability.}{Marketplace de doações em .NET com PostgreSQL, voltado a economia circular e escalabilidade.}
\end{itemize}

\section*{\tr{Education}{Formação Acadêmica}}

\entry{\tr{Technologist in Systems Analysis and Development}{Tecnólogo em Análise e Desenvolvimento de Sistemas}}{2025}{Unicesumar}{Londrina - PR}

\section*{\tr{Certifications}{Certificações}}

\entry{Foundational C\# with Microsoft}{06/2025}{\tr{.NET and C\# fundamentals}{Fundamentos de .NET e C\#}}{\href{https://freecodecamp.org/certification/h0wzy/foundational-c-sharp-with-microsoft}{\faLink \space freeCodeCamp}}
\vspace{-0.4em}
\begin{itemize}[leftmargin=*, itemsep=2pt]
\setlength\itemsep{0em}
\item \tr{C\# syntax and control structures, OOP, collections, exception handling and .NET best practices.}{Sintaxe e estruturas de C\#, POO, coleções, tratamento de exceções e boas práticas em .NET.}
\end{itemize}

\end{document}
)
%  Macro de tradução:  \tr{<texto EN>}{<texto PT>}
% =====================================================================

\documentclass[a4paper,10pt]{article}
\usepackage[margin=0.5in,nofoot]{geometry}
\usepackage{fontawesome5}
\usepackage{xcolor}

\definecolor{accent}{HTML}{1F4E79}

% ---- Fonte (XeLaTeX/LuaLaTeX) ----
\usepackage{fontspec}
\setmainfont{Montserrat}[
    Path = ./fonts/,
    Extension = .ttf,
    UprightFont = *-Regular,
    BoldFont = *-Bold,
    ItalicFont = *-Italic
]

% ---- Idioma / hifenização ----
\usepackage{polyglossia}
\ifdefined\CVPT
  \setdefaultlanguage{portuguese}\setotherlanguage{english}
\else
  \setdefaultlanguage{english}\setotherlanguage{portuguese}
\fi

% ---- Macro de tradução: \tr{EN}{PT} (default = EN) ----
\newcommand{\tr}[2]{\ifdefined\CVPT #2\else #1\fi}

\usepackage{titlesec}
\usepackage{enumitem}

\usepackage{hyperref}
\hypersetup{
    colorlinks=true,
    linkcolor=accent,
    filecolor=accent,
    urlcolor=accent,
    citecolor=accent
}

\titleformat{\section}{\large\bfseries}{\thesection}{1em}{}[\titlerule]
\titlespacing*{\section}{0pt}{*1}{*1}

\newcommand{\entry}[4]{
  \noindent\textbf{#1} \hfill #2 \\
  \noindent\textit{#3} \hfill \textit{#4} \\
  \vspace{2pt}
}

\newcommand{\project}[2]{
  \noindent\textbf{#1} \hfill #2 \\
  \vspace{2pt}
}

\begin{document}
\small
\setlength{\parskip}{4pt}
\setlength{\parindent}{0pt}

\pagenumbering{gobble}

\noindent
\begin{minipage}[t]{0.5\textwidth}
{\Large \textbf{Marcos Junior Bueno Selzler}}

\vspace{0.4em}

\end{minipage}%
\begin{minipage}[t]{0.5\textwidth}
\raggedleft

\tr{Brazil}{Brasil} (GMT-3)

\href{tel:+5543988431018}{\faPhone \space (43) 98843-1018} \\
\href{mailto:h0wzymarcos@gmail.com}{\faEnvelope \space h0wzymarcos@gmail.com} \\
\href{https://www.linkedin.com/in/marcosh0wzy/}{\faLinkedin \space marcosh0wzy} \\
\href{https://github.com/H0wZy}{\faGithub \space H0wZy} \\
\href{https://linktr.ee/h0wzymarcos}{\faLink \space \tr{Portfolio}{Portfólio}}

\end{minipage}

\vspace{1em}

\begin{center}
    \textbf{SOFTWARE ENGINEER \textbar\ BACKEND \& FULL STACK}
\end{center}

\vspace{0.4em}

\noindent
\tr{%
Software Engineer at TCS (Copel account), focused on backend with Java/JSF and .NET/Golang. Two-time TCS AI Friday champion for solutions built with LLMs and predictive models. Experienced in modernizing legacy systems to the cloud (GCP/AWS) and building REST APIs with Clean Code, SOLID and testing.%
}{%
Engenheiro de Software na TCS (conta Copel), com foco em backend Java/JSF e .NET/Golang. Bicampeão do TCS AI Friday por soluções com LLMs e modelos preditivos. Experiência em modernização de sistemas legados para cloud (GCP/AWS) e construção de APIs REST com Clean Code, SOLID e testes.%
}

\section*{\tr{Professional Experience}{Experiência Profissional}}

\entry{Tata Consultancy Services \textbar\ Copel}{\faCalendar \space 05/2025 -- \tr{Present}{Presente}}{Software Engineer}{Londrina - PR}
\vspace{-0.4em}
\begin{itemize}[leftmargin=*, itemsep=2pt]
\setlength\itemsep{0em}
\item \tr{I develop and maintain Java 8 / JSF applications on WildFly/JBoss servers.}{Desenvolvo e mantenho aplicações Java 8 com JSF em servidores WildFly/JBoss.}
\item \tr{Optimized critical SQL queries, improving performance and stability.}{Otimizei queries SQL críticas, melhorando performance e estabilidade.}
\item \tr{Migrated applications from JBoss 7 to WildFly 23.}{Migrei aplicações de JBoss 7 para WildFly 23.}
\item \tr{Modernize on-premise legacy systems to GCP, with integrations via Denodo.}{Modernizo sistemas legados on-premise para o GCP, com integrações via Denodo.}
\item \tr{Build backend services in Golang.}{Desenvolvo serviços de backend em Golang.}
\item \tr{Work across cloud environments (GCP and AWS).}{Atuo em ambientes cloud (GCP e AWS).}
\end{itemize}

\entry{Tata Consultancy Services \textbar\ Bootcamp 2x2y}{\faCalendar \space 05/2025 -- 11/2025}{Bootcamper Software Engineer}{Londrina - PR}
\vspace{-0.4em}
\begin{itemize}[leftmargin=*, itemsep=2pt]
\setlength\itemsep{0em}
\item \tr{Intensive full-stack software engineering training at TCS.}{Formação intensiva em engenharia de software full stack pela TCS.}
\item \tr{Built hands-on projects applying layered architecture and Clean Code.}{Desenvolvi projetos práticos aplicando arquitetura em camadas e Clean Code.}
\end{itemize}

\section*{\tr{Awards \& Highlights}{Prêmios e Destaques}}

\entry{\tr{Two-time TCS AI Friday Champion}{Bicampeão TCS AI Friday}}{05/2026}{\tr{Pipeverse Project}{Projeto Pipeverse}}{\href{https://www.linkedin.com/posts/marcosh0wzy_tcs-tcsai-aifriday-activity-7463971325016006657-WzoZ}{\faLink \space LinkedIn}}
\vspace{-0.4em}
\begin{itemize}[leftmargin=*, itemsep=2pt]
\setlength\itemsep{0em}
\item \tr{Predictive system to estimate the probability of failure in AI projects (data quality, problem definition, adoption and monitoring).}{Sistema preditivo para estimar a probabilidade de falha em projetos de IA (qualidade de dados, definição do problema, adoção e monitoramento).}
\end{itemize}

\entry{\tr{TCS AI Friday Champion}{Campeão TCS AI Friday}}{02/2026}{\tr{Energy Saver Project}{Projeto Energy Saver}}{\href{https://www.linkedin.com/posts/marcosh0wzy_ai-llm-dataanalysis-activity-7439643718057385985--a7G}{\faLink \space LinkedIn}}
\vspace{-0.4em}
\begin{itemize}[leftmargin=*, itemsep=2pt]
\setlength\itemsep{0em}
\item \tr{LLM-based solution for energy-consumption analysis from CSV data: anomaly detection and automated recommendations to reduce waste.}{Solução com LLM para análise de consumo energético a partir de CSV: detecção de anomalias e recomendações automatizadas para redução de desperdício.}
\end{itemize}

\section*{\tr{Skills}{Habilidades}}

\begin{itemize}[leftmargin=*, itemsep=2pt]
\setlength\itemsep{0em}
\item \textbf{Tech Stack}: .NET, C\#, Golang, Java, Spring Boot, React, Next.js, TypeScript, SQL/NoSQL, Docker
\item \textbf{Cloud \& DevOps}: GCP (Cloud Run, Cloud SQL, GCS), AWS, Terraform (IaC), Docker, CI/CD, Git/GitHub
\item \textbf{\tr{AI \& LLMs}{IA \& LLMs}}: \tr{data analysis with LLMs, predictive models, prompt engineering}{análise de dados com LLMs, modelos preditivos, prompt engineering}
\item \textbf{\tr{Concepts}{Conceitos}}: \tr{Clean Code, SOLID, DDD, OOP, REST APIs, Observability, Unit Testing (xUnit, Moq)}{Clean Code, SOLID, DDD, POO, APIs REST, Observabilidade, Testes Unitários (xUnit, Moq)}
\item \textbf{\tr{Languages}{Idiomas}}: \tr{English (Intermediate), Portuguese (Native)}{Inglês (Intermediário), Português (Nativo)}
\end{itemize}

\section*{\tr{Projects}{Projetos}}

\project{Telas Paraná --- \tr{Institutional Website}{Site Institucional}}{\href{https://frontend-4mxyivrfva-r.a.run.app/}{\faLink \space \tr{Live}{Produção}} \textbar\ \href{https://telasearamesparana.com.br}{\tr{previous site}{site anterior}}}
\vspace{-0.4em}
\begin{itemize}[leftmargin=*, itemsep=2pt]
\setlength\itemsep{0em}
\item \tr{Full-stack redesign of the company's institutional website, replacing the legacy site.}{Redesign full-stack do site institucional da empresa, substituindo o site legado.}
\item \tr{React / Next.js front-end and two .NET/C\# microservices (Users.Api, Leads.Api), all deployed on GCP Cloud Run.}{Front-end em React / Next.js e dois microsserviços .NET/C\# (Users.Api, Leads.Api), todos implantados no GCP Cloud Run.}
\item \tr{Cloud SQL database, Cloud Storage (GCS) for media, and infrastructure provisioned as code with Terraform.}{Banco no Cloud SQL, mídia no Cloud Storage (GCS) e infraestrutura provisionada como código com Terraform.}
\end{itemize}

\project{C\# User API}{\href{https://github.com/H0wZy/UserApi}{\faLink \space GitHub}}
\vspace{-0.4em}
\begin{itemize}[leftmargin=*, itemsep=2pt]
\setlength\itemsep{0em}
\item \tr{REST API in C\#/.NET with Entity Framework and PostgreSQL, layered architecture and user CRUD.}{API REST em C\#/.NET com Entity Framework e PostgreSQL, arquitetura em camadas e CRUD de usuários.}
\item \tr{JWT authentication, authorization control and test coverage with xUnit and Moq.}{Autenticação JWT, controle de autorização e cobertura de testes com xUnit e Moq.}
\end{itemize}

\project{Golang User API}{\href{https://github.com/H0wZy/user-api}{\faLink \space GitHub}}
\vspace{-0.4em}
\begin{itemize}[leftmargin=*, itemsep=2pt]
\setlength\itemsep{0em}
\item \tr{REST API in Golang with Gin and Gorm: user CRUD, authentication and authorization.}{API REST em Golang com Gin e Gorm: CRUD de usuários, autenticação e autorização.}
\item \tr{Backend structured with a focus on performance and code organization.}{Backend estruturado com foco em performance e organização de código.}
\end{itemize}

\project{Web EasyDonate}{\href{https://github.com/H0wZy/web.easydonate}{\faLink \space GitHub}}
\vspace{-0.4em}
\begin{itemize}[leftmargin=*, itemsep=2pt]
\setlength\itemsep{0em}
\item \tr{Donation marketplace in .NET with PostgreSQL, oriented to circular economy and scalability.}{Marketplace de doações em .NET com PostgreSQL, voltado a economia circular e escalabilidade.}
\end{itemize}

\section*{\tr{Education}{Formação Acadêmica}}

\entry{\tr{Technologist in Systems Analysis and Development}{Tecnólogo em Análise e Desenvolvimento de Sistemas}}{2025}{Unicesumar}{Londrina - PR}

\section*{\tr{Certifications}{Certificações}}

\entry{Foundational C\# with Microsoft}{06/2025}{\tr{.NET and C\# fundamentals}{Fundamentos de .NET e C\#}}{\href{https://freecodecamp.org/certification/h0wzy/foundational-c-sharp-with-microsoft}{\faLink \space freeCodeCamp}}
\vspace{-0.4em}
\begin{itemize}[leftmargin=*, itemsep=2pt]
\setlength\itemsep{0em}
\item \tr{C\# syntax and control structures, OOP, collections, exception handling and .NET best practices.}{Sintaxe e estruturas de C\#, POO, coleções, tratamento de exceções e boas práticas em .NET.}
\end{itemize}

\end{document}
)
%  Macro de tradução:  \tr{<texto EN>}{<texto PT>}
% =====================================================================

\documentclass[a4paper,10pt]{article}
\usepackage[margin=0.5in,nofoot]{geometry}
\usepackage{fontawesome5}
\usepackage{xcolor}

\definecolor{accent}{HTML}{1F4E79}

% ---- Fonte (XeLaTeX/LuaLaTeX) ----
\usepackage{fontspec}
\setmainfont{Montserrat}[
    Path = ./fonts/,
    Extension = .ttf,
    UprightFont = *-Regular,
    BoldFont = *-Bold,
    ItalicFont = *-Italic
]

% ---- Idioma / hifenização ----
\usepackage{polyglossia}
\ifdefined\CVPT
  \setdefaultlanguage{portuguese}\setotherlanguage{english}
\else
  \setdefaultlanguage{english}\setotherlanguage{portuguese}
\fi

% ---- Macro de tradução: \tr{EN}{PT} (default = EN) ----
\newcommand{\tr}[2]{\ifdefined\CVPT #2\else #1\fi}

\usepackage{titlesec}
\usepackage{enumitem}

\usepackage{hyperref}
\hypersetup{
    colorlinks=true,
    linkcolor=accent,
    filecolor=accent,
    urlcolor=accent,
    citecolor=accent
}

\titleformat{\section}{\large\bfseries}{\thesection}{1em}{}[\titlerule]
\titlespacing*{\section}{0pt}{*1}{*1}

\newcommand{\entry}[4]{
  \noindent\textbf{#1} \hfill #2 \\
  \noindent\textit{#3} \hfill \textit{#4} \\
  \vspace{2pt}
}

\newcommand{\project}[2]{
  \noindent\textbf{#1} \hfill #2 \\
  \vspace{2pt}
}

\begin{document}
\small
\setlength{\parskip}{4pt}
\setlength{\parindent}{0pt}

\pagenumbering{gobble}

\noindent
\begin{minipage}[t]{0.5\textwidth}
{\Large \textbf{Marcos Junior Bueno Selzler}}

\vspace{0.4em}

\end{minipage}%
\begin{minipage}[t]{0.5\textwidth}
\raggedleft

\tr{Brazil}{Brasil} (GMT-3)

\href{tel:+5543988431018}{\faPhone \space (43) 98843-1018} \\
\href{mailto:h0wzymarcos@gmail.com}{\faEnvelope \space h0wzymarcos@gmail.com} \\
\href{https://www.linkedin.com/in/marcosh0wzy/}{\faLinkedin \space marcosh0wzy} \\
\href{https://github.com/H0wZy}{\faGithub \space H0wZy} \\
\href{https://linktr.ee/h0wzymarcos}{\faLink \space \tr{Portfolio}{Portfólio}}

\end{minipage}

\vspace{1em}

\begin{center}
    \textbf{SOFTWARE ENGINEER \textbar\ BACKEND \& FULL STACK}
\end{center}

\vspace{0.4em}

\noindent
\tr{%
Software Engineer at TCS (Copel account), focused on backend with Java/JSF and .NET/Golang. Two-time TCS AI Friday champion for solutions built with LLMs and predictive models. Experienced in modernizing legacy systems to the cloud (GCP/AWS) and building REST APIs with Clean Code, SOLID and testing.%
}{%
Engenheiro de Software na TCS (conta Copel), com foco em backend Java/JSF e .NET/Golang. Bicampeão do TCS AI Friday por soluções com LLMs e modelos preditivos. Experiência em modernização de sistemas legados para cloud (GCP/AWS) e construção de APIs REST com Clean Code, SOLID e testes.%
}

\section*{\tr{Professional Experience}{Experiência Profissional}}

\entry{Tata Consultancy Services \textbar\ Copel}{\faCalendar \space 05/2025 -- \tr{Present}{Presente}}{Software Engineer}{Londrina - PR}
\vspace{-0.4em}
\begin{itemize}[leftmargin=*, itemsep=2pt]
\setlength\itemsep{0em}
\item \tr{I develop and maintain Java 8 / JSF applications on WildFly/JBoss servers.}{Desenvolvo e mantenho aplicações Java 8 com JSF em servidores WildFly/JBoss.}
\item \tr{Optimized critical SQL queries, improving performance and stability.}{Otimizei queries SQL críticas, melhorando performance e estabilidade.}
\item \tr{Migrated applications from JBoss 7 to WildFly 23.}{Migrei aplicações de JBoss 7 para WildFly 23.}
\item \tr{Modernize on-premise legacy systems to GCP, with integrations via Denodo.}{Modernizo sistemas legados on-premise para o GCP, com integrações via Denodo.}
\item \tr{Build backend services in Golang.}{Desenvolvo serviços de backend em Golang.}
\item \tr{Work across cloud environments (GCP and AWS).}{Atuo em ambientes cloud (GCP e AWS).}
\end{itemize}

\entry{Tata Consultancy Services \textbar\ Bootcamp 2x2y}{\faCalendar \space 05/2025 -- 11/2025}{Bootcamper Software Engineer}{Londrina - PR}
\vspace{-0.4em}
\begin{itemize}[leftmargin=*, itemsep=2pt]
\setlength\itemsep{0em}
\item \tr{Intensive full-stack software engineering training at TCS.}{Formação intensiva em engenharia de software full stack pela TCS.}
\item \tr{Built hands-on projects applying layered architecture and Clean Code.}{Desenvolvi projetos práticos aplicando arquitetura em camadas e Clean Code.}
\end{itemize}

\section*{\tr{Awards \& Highlights}{Prêmios e Destaques}}

\entry{\tr{Two-time TCS AI Friday Champion}{Bicampeão TCS AI Friday}}{05/2026}{\tr{Pipeverse Project}{Projeto Pipeverse}}{\href{https://www.linkedin.com/posts/marcosh0wzy_tcs-tcsai-aifriday-activity-7463971325016006657-WzoZ}{\faLink \space LinkedIn}}
\vspace{-0.4em}
\begin{itemize}[leftmargin=*, itemsep=2pt]
\setlength\itemsep{0em}
\item \tr{Predictive system to estimate the probability of failure in AI projects (data quality, problem definition, adoption and monitoring).}{Sistema preditivo para estimar a probabilidade de falha em projetos de IA (qualidade de dados, definição do problema, adoção e monitoramento).}
\end{itemize}

\entry{\tr{TCS AI Friday Champion}{Campeão TCS AI Friday}}{02/2026}{\tr{Energy Saver Project}{Projeto Energy Saver}}{\href{https://www.linkedin.com/posts/marcosh0wzy_ai-llm-dataanalysis-activity-7439643718057385985--a7G}{\faLink \space LinkedIn}}
\vspace{-0.4em}
\begin{itemize}[leftmargin=*, itemsep=2pt]
\setlength\itemsep{0em}
\item \tr{LLM-based solution for energy-consumption analysis from CSV data: anomaly detection and automated recommendations to reduce waste.}{Solução com LLM para análise de consumo energético a partir de CSV: detecção de anomalias e recomendações automatizadas para redução de desperdício.}
\end{itemize}

\section*{\tr{Skills}{Habilidades}}

\begin{itemize}[leftmargin=*, itemsep=2pt]
\setlength\itemsep{0em}
\item \textbf{Tech Stack}: .NET, C\#, Golang, Java, Spring Boot, React, Next.js, TypeScript, SQL/NoSQL, Docker
\item \textbf{Cloud \& DevOps}: GCP (Cloud Run, Cloud SQL, GCS), AWS, Terraform (IaC), Docker, CI/CD, Git/GitHub
\item \textbf{\tr{AI \& LLMs}{IA \& LLMs}}: \tr{data analysis with LLMs, predictive models, prompt engineering}{análise de dados com LLMs, modelos preditivos, prompt engineering}
\item \textbf{\tr{Concepts}{Conceitos}}: \tr{Clean Code, SOLID, DDD, OOP, REST APIs, Observability, Unit Testing (xUnit, Moq)}{Clean Code, SOLID, DDD, POO, APIs REST, Observabilidade, Testes Unitários (xUnit, Moq)}
\item \textbf{\tr{Languages}{Idiomas}}: \tr{English (Intermediate), Portuguese (Native)}{Inglês (Intermediário), Português (Nativo)}
\end{itemize}

\section*{\tr{Projects}{Projetos}}

\project{Telas Paraná --- \tr{Institutional Website}{Site Institucional}}{\href{https://frontend-4mxyivrfva-r.a.run.app/}{\faLink \space \tr{Live}{Produção}} \textbar\ \href{https://telasearamesparana.com.br}{\tr{previous site}{site anterior}}}
\vspace{-0.4em}
\begin{itemize}[leftmargin=*, itemsep=2pt]
\setlength\itemsep{0em}
\item \tr{Full-stack redesign of the company's institutional website, replacing the legacy site.}{Redesign full-stack do site institucional da empresa, substituindo o site legado.}
\item \tr{React / Next.js front-end and two .NET/C\# microservices (Users.Api, Leads.Api), all deployed on GCP Cloud Run.}{Front-end em React / Next.js e dois microsserviços .NET/C\# (Users.Api, Leads.Api), todos implantados no GCP Cloud Run.}
\item \tr{Cloud SQL database, Cloud Storage (GCS) for media, and infrastructure provisioned as code with Terraform.}{Banco no Cloud SQL, mídia no Cloud Storage (GCS) e infraestrutura provisionada como código com Terraform.}
\end{itemize}

\project{C\# User API}{\href{https://github.com/H0wZy/UserApi}{\faLink \space GitHub}}
\vspace{-0.4em}
\begin{itemize}[leftmargin=*, itemsep=2pt]
\setlength\itemsep{0em}
\item \tr{REST API in C\#/.NET with Entity Framework and PostgreSQL, layered architecture and user CRUD.}{API REST em C\#/.NET com Entity Framework e PostgreSQL, arquitetura em camadas e CRUD de usuários.}
\item \tr{JWT authentication, authorization control and test coverage with xUnit and Moq.}{Autenticação JWT, controle de autorização e cobertura de testes com xUnit e Moq.}
\end{itemize}

\project{Golang User API}{\href{https://github.com/H0wZy/user-api}{\faLink \space GitHub}}
\vspace{-0.4em}
\begin{itemize}[leftmargin=*, itemsep=2pt]
\setlength\itemsep{0em}
\item \tr{REST API in Golang with Gin and Gorm: user CRUD, authentication and authorization.}{API REST em Golang com Gin e Gorm: CRUD de usuários, autenticação e autorização.}
\item \tr{Backend structured with a focus on performance and code organization.}{Backend estruturado com foco em performance e organização de código.}
\end{itemize}

\project{Web EasyDonate}{\href{https://github.com/H0wZy/web.easydonate}{\faLink \space GitHub}}
\vspace{-0.4em}
\begin{itemize}[leftmargin=*, itemsep=2pt]
\setlength\itemsep{0em}
\item \tr{Donation marketplace in .NET with PostgreSQL, oriented to circular economy and scalability.}{Marketplace de doações em .NET com PostgreSQL, voltado a economia circular e escalabilidade.}
\end{itemize}

\section*{\tr{Education}{Formação Acadêmica}}

\entry{\tr{Technologist in Systems Analysis and Development}{Tecnólogo em Análise e Desenvolvimento de Sistemas}}{2025}{Unicesumar}{Londrina - PR}

\section*{\tr{Certifications}{Certificações}}

\entry{Foundational C\# with Microsoft}{06/2025}{\tr{.NET and C\# fundamentals}{Fundamentos de .NET e C\#}}{\href{https://freecodecamp.org/certification/h0wzy/foundational-c-sharp-with-microsoft}{\faLink \space freeCodeCamp}}
\vspace{-0.4em}
\begin{itemize}[leftmargin=*, itemsep=2pt]
\setlength\itemsep{0em}
\item \tr{C\# syntax and control structures, OOP, collections, exception handling and .NET best practices.}{Sintaxe e estruturas de C\#, POO, coleções, tratamento de exceções e boas práticas em .NET.}
\end{itemize}

\end{document}
)
%  Macro de tradução:  \tr{<texto EN>}{<texto PT>}
% =====================================================================

\documentclass[a4paper,10pt]{article}
\usepackage[margin=0.5in,nofoot]{geometry}
\usepackage{fontawesome5}
\usepackage{xcolor}

\definecolor{accent}{HTML}{1F4E79}

% ---- Fonte (XeLaTeX/LuaLaTeX) ----
\usepackage{fontspec}
\setmainfont{Montserrat}[
    Path = ./fonts/,
    Extension = .ttf,
    UprightFont = *-Regular,
    BoldFont = *-Bold,
    ItalicFont = *-Italic
]

% ---- Idioma / hifenização ----
\usepackage{polyglossia}
\ifdefined\CVPT
  \setdefaultlanguage{portuguese}\setotherlanguage{english}
\else
  \setdefaultlanguage{english}\setotherlanguage{portuguese}
\fi

% ---- Macro de tradução: \tr{EN}{PT} (default = EN) ----
\newcommand{\tr}[2]{\ifdefined\CVPT #2\else #1\fi}

\usepackage{titlesec}
\usepackage{enumitem}

\usepackage{hyperref}
\hypersetup{
    colorlinks=true,
    linkcolor=accent,
    filecolor=accent,
    urlcolor=accent,
    citecolor=accent
}

\titleformat{\section}{\large\bfseries}{\thesection}{1em}{}[\titlerule]
\titlespacing*{\section}{0pt}{*1}{*1}

\newcommand{\entry}[4]{
  \noindent\textbf{#1} \hfill #2 \\
  \noindent\textit{#3} \hfill \textit{#4} \\
  \vspace{2pt}
}

\newcommand{\project}[2]{
  \noindent\textbf{#1} \hfill #2 \\
  \vspace{2pt}
}

\begin{document}
\small
\setlength{\parskip}{4pt}
\setlength{\parindent}{0pt}

\pagenumbering{gobble}

\noindent
\begin{minipage}[t]{0.5\textwidth}
{\Large \textbf{Marcos Junior Bueno Selzler}}

\vspace{0.4em}

\end{minipage}%
\begin{minipage}[t]{0.5\textwidth}
\raggedleft

\tr{Brazil}{Brasil} (GMT-3)

\href{tel:+5543988431018}{\faPhone \space (43) 98843-1018} \\
\href{mailto:h0wzymarcos@gmail.com}{\faEnvelope \space h0wzymarcos@gmail.com} \\
\href{https://www.linkedin.com/in/marcosh0wzy/}{\faLinkedin \space marcosh0wzy} \\
\href{https://github.com/H0wZy}{\faGithub \space H0wZy} \\
\href{https://linktr.ee/h0wzymarcos}{\faLink \space \tr{Portfolio}{Portfólio}}

\end{minipage}

\vspace{1em}

\begin{center}
    \textbf{SOFTWARE ENGINEER \textbar\ BACKEND \& FULL STACK}
\end{center}

\vspace{0.4em}

\noindent
\tr{%
Software Engineer at TCS (Copel account), focused on backend with Java/JSF and .NET/Golang. Two-time TCS AI Friday champion for solutions built with LLMs and predictive models. Experienced in modernizing legacy systems to the cloud (GCP/AWS) and building REST APIs with Clean Code, SOLID and testing.%
}{%
Engenheiro de Software na TCS (conta Copel), com foco em backend Java/JSF e .NET/Golang. Bicampeão do TCS AI Friday por soluções com LLMs e modelos preditivos. Experiência em modernização de sistemas legados para cloud (GCP/AWS) e construção de APIs REST com Clean Code, SOLID e testes.%
}

\section*{\tr{Professional Experience}{Experiência Profissional}}

\entry{Tata Consultancy Services \textbar\ Copel}{\faCalendar \space 05/2025 -- \tr{Present}{Presente}}{Software Engineer}{Londrina - PR}
\vspace{-0.4em}
\begin{itemize}[leftmargin=*, itemsep=2pt]
\setlength\itemsep{0em}
\item \tr{I develop and maintain Java 8 / JSF applications on WildFly/JBoss servers.}{Desenvolvo e mantenho aplicações Java 8 com JSF em servidores WildFly/JBoss.}
\item \tr{Optimized critical SQL queries, improving performance and stability.}{Otimizei queries SQL críticas, melhorando performance e estabilidade.}
\item \tr{Migrated applications from JBoss 7 to WildFly 23.}{Migrei aplicações de JBoss 7 para WildFly 23.}
\item \tr{Modernize on-premise legacy systems to GCP, with integrations via Denodo.}{Modernizo sistemas legados on-premise para o GCP, com integrações via Denodo.}
\item \tr{Build backend services in Golang.}{Desenvolvo serviços de backend em Golang.}
\item \tr{Work across cloud environments (GCP and AWS).}{Atuo em ambientes cloud (GCP e AWS).}
\end{itemize}

\entry{Tata Consultancy Services \textbar\ Bootcamp 2x2y}{\faCalendar \space 05/2025 -- 11/2025}{Bootcamper Software Engineer}{Londrina - PR}
\vspace{-0.4em}
\begin{itemize}[leftmargin=*, itemsep=2pt]
\setlength\itemsep{0em}
\item \tr{Intensive full-stack software engineering training at TCS.}{Formação intensiva em engenharia de software full stack pela TCS.}
\item \tr{Built hands-on projects applying layered architecture and Clean Code.}{Desenvolvi projetos práticos aplicando arquitetura em camadas e Clean Code.}
\end{itemize}

\section*{\tr{Awards \& Highlights}{Prêmios e Destaques}}

\entry{\tr{Two-time TCS AI Friday Champion}{Bicampeão TCS AI Friday}}{05/2026}{\tr{Pipeverse Project}{Projeto Pipeverse}}{\href{https://www.linkedin.com/posts/marcosh0wzy_tcs-tcsai-aifriday-activity-7463971325016006657-WzoZ}{\faLink \space LinkedIn}}
\vspace{-0.4em}
\begin{itemize}[leftmargin=*, itemsep=2pt]
\setlength\itemsep{0em}
\item \tr{Predictive system to estimate the probability of failure in AI projects (data quality, problem definition, adoption and monitoring).}{Sistema preditivo para estimar a probabilidade de falha em projetos de IA (qualidade de dados, definição do problema, adoção e monitoramento).}
\end{itemize}

\entry{\tr{TCS AI Friday Champion}{Campeão TCS AI Friday}}{02/2026}{\tr{Energy Saver Project}{Projeto Energy Saver}}{\href{https://www.linkedin.com/posts/marcosh0wzy_ai-llm-dataanalysis-activity-7439643718057385985--a7G}{\faLink \space LinkedIn}}
\vspace{-0.4em}
\begin{itemize}[leftmargin=*, itemsep=2pt]
\setlength\itemsep{0em}
\item \tr{LLM-based solution for energy-consumption analysis from CSV data: anomaly detection and automated recommendations to reduce waste.}{Solução com LLM para análise de consumo energético a partir de CSV: detecção de anomalias e recomendações automatizadas para redução de desperdício.}
\end{itemize}

\section*{\tr{Skills}{Habilidades}}

\begin{itemize}[leftmargin=*, itemsep=2pt]
\setlength\itemsep{0em}
\item \textbf{Tech Stack}: .NET, C\#, Golang, Java, Spring Boot, React, Next.js, TypeScript, SQL/NoSQL, Docker
\item \textbf{Cloud \& DevOps}: GCP (Cloud Run, Cloud SQL, GCS), AWS, Terraform (IaC), Docker, CI/CD, Git/GitHub
\item \textbf{\tr{AI \& LLMs}{IA \& LLMs}}: \tr{data analysis with LLMs, predictive models, prompt engineering}{análise de dados com LLMs, modelos preditivos, prompt engineering}
\item \textbf{\tr{Concepts}{Conceitos}}: \tr{Clean Code, SOLID, DDD, OOP, REST APIs, Observability, Unit Testing (xUnit, Moq)}{Clean Code, SOLID, DDD, POO, APIs REST, Observabilidade, Testes Unitários (xUnit, Moq)}
\item \textbf{\tr{Languages}{Idiomas}}: \tr{English (Intermediate), Portuguese (Native)}{Inglês (Intermediário), Português (Nativo)}
\end{itemize}

\section*{\tr{Projects}{Projetos}}

\project{Telas Paraná --- \tr{Institutional Website}{Site Institucional}}{\href{https://frontend-4mxyivrfva-r.a.run.app/}{\faLink \space \tr{Live}{Produção}} \textbar\ \href{https://telasearamesparana.com.br}{\tr{previous site}{site anterior}}}
\vspace{-0.4em}
\begin{itemize}[leftmargin=*, itemsep=2pt]
\setlength\itemsep{0em}
\item \tr{Full-stack redesign of the company's institutional website, replacing the legacy site.}{Redesign full-stack do site institucional da empresa, substituindo o site legado.}
\item \tr{React / Next.js front-end and two .NET/C\# microservices (Users.Api, Leads.Api), all deployed on GCP Cloud Run.}{Front-end em React / Next.js e dois microsserviços .NET/C\# (Users.Api, Leads.Api), todos implantados no GCP Cloud Run.}
\item \tr{Cloud SQL database, Cloud Storage (GCS) for media, and infrastructure provisioned as code with Terraform.}{Banco no Cloud SQL, mídia no Cloud Storage (GCS) e infraestrutura provisionada como código com Terraform.}
\end{itemize}

\project{C\# User API}{\href{https://github.com/H0wZy/UserApi}{\faLink \space GitHub}}
\vspace{-0.4em}
\begin{itemize}[leftmargin=*, itemsep=2pt]
\setlength\itemsep{0em}
\item \tr{REST API in C\#/.NET with Entity Framework and PostgreSQL, layered architecture and user CRUD.}{API REST em C\#/.NET com Entity Framework e PostgreSQL, arquitetura em camadas e CRUD de usuários.}
\item \tr{JWT authentication, authorization control and test coverage with xUnit and Moq.}{Autenticação JWT, controle de autorização e cobertura de testes com xUnit e Moq.}
\end{itemize}

\project{Golang User API}{\href{https://github.com/H0wZy/user-api}{\faLink \space GitHub}}
\vspace{-0.4em}
\begin{itemize}[leftmargin=*, itemsep=2pt]
\setlength\itemsep{0em}
\item \tr{REST API in Golang with Gin and Gorm: user CRUD, authentication and authorization.}{API REST em Golang com Gin e Gorm: CRUD de usuários, autenticação e autorização.}
\item \tr{Backend structured with a focus on performance and code organization.}{Backend estruturado com foco em performance e organização de código.}
\end{itemize}

\project{Web EasyDonate}{\href{https://github.com/H0wZy/web.easydonate}{\faLink \space GitHub}}
\vspace{-0.4em}
\begin{itemize}[leftmargin=*, itemsep=2pt]
\setlength\itemsep{0em}
\item \tr{Donation marketplace in .NET with PostgreSQL, oriented to circular economy and scalability.}{Marketplace de doações em .NET com PostgreSQL, voltado a economia circular e escalabilidade.}
\end{itemize}

\section*{\tr{Education}{Formação Acadêmica}}

\entry{\tr{Technologist in Systems Analysis and Development}{Tecnólogo em Análise e Desenvolvimento de Sistemas}}{2025}{Unicesumar}{Londrina - PR}

\section*{\tr{Certifications}{Certificações}}

\entry{Foundational C\# with Microsoft}{06/2025}{\tr{.NET and C\# fundamentals}{Fundamentos de .NET e C\#}}{\href{https://freecodecamp.org/certification/h0wzy/foundational-c-sharp-with-microsoft}{\faLink \space freeCodeCamp}}
\vspace{-0.4em}
\begin{itemize}[leftmargin=*, itemsep=2pt]
\setlength\itemsep{0em}
\item \tr{C\# syntax and control structures, OOP, collections, exception handling and .NET best practices.}{Sintaxe e estruturas de C\#, POO, coleções, tratamento de exceções e boas práticas em .NET.}
\end{itemize}

\end{document}
